\begin{table*}[t]
\centering
\caption{Curation failure modes and metadata fixes.}
\label{tab:failure-examples}
\small
\begin{tabular}{p{0.22\linewidth}p{0.34\linewidth}p{0.34\linewidth}}
\toprule
Failure mode & Curation problem & Metadata fix \\
\midrule
Artifact role ambiguity & A single repository can expose base-model, merge, quantization, or deployment signals at the same time. & Store primary artifact role separately from packaging overlays. \\
Parent-only provenance & Family labels help discovery but can mislead if treated as exact upstream claims. & Use typed upstream relations such as base, teacher, adapter target, merge input, quantization source, and baseline. \\
Teacher/reference confusion & Teacher or evaluator mentions can be mistaken for the artifact's own family. & Separate target artifact, teacher model, evaluator, and baseline fields. \\
False positive family link & Text mentions alone overstate provenance precision. & Attach provenance confidence and evidence type to family labels. \\
Stated release without reusable asset & Binary open/not-open metadata hides large differences in release evidence and preservation triage. & Expose separate release fields for code, weights/adapters, data/evaluation assets, scripts, and complete recipes. \\
No bidirectional paper-repository link & Curators cannot reliably connect scholarly claims to artifact records. & Add DOI/arXiv-to-repository and repository-to-paper relation fields. \\
\bottomrule
\end{tabular}
\end{table*}
